% =====================================================================
%  CHAPTER 3: 儿童少年身体发育（对应大纲第四章）
%  参考往届笔记：第三章 儿童少年身体发育
% =====================================================================
\chapter{儿童少年身体发育}
\setcounter{chapter}{3}
\chapterintro{
\textbf{【掌握】}儿童少年体格发育（体格阶段性变化、身体比例变化、体型变化）；儿童少年体能发育（健康相关体能、运动技能相关体能、体能发育特点、体力活动）。\par
\textbf{【熟悉】}儿童少年体成分发育（体成分模型、体成分发育的性别特征）。
}

% ----- 5.1 体格发育 -----
\section{儿童少年体格发育}

\subsection{体格阶段性变化}

\begin{keybox}
\textbf{体格发育（Physical Growth）：}人体外部形态、身体比例和体型等方面随年龄发生的变化。

\textbf{一、体格阶段性变化：}身高$\to$身体线性生长（Linear Growth），体重$\to$重量。

\textbf{1. 第一生长突增期：}胎儿$\to$2岁。
\begin{itemize}[leftmargin=*]
    \item \textbf{身高：}孕中期$\to$1岁末；孕中期（4-6个月）增长量占胎儿期总增长量的50\%；出生后第1年增长约25cm。
    \item \textbf{体重：}孕晚期$\to$1岁末；孕晚期（7-9个月）增长量占胎儿期总增长量的70\%。
\end{itemize}

\textbf{2. 相对稳定期：}2岁$\to$青春期发育前。身高：$\uparrow$5-7cm/年，体重：$\uparrow$2-3kg/年。

\textbf{3. 第二生长突增期/青春期生长突增（Adolescent Growth Spurt）：}青春期。
\begin{itemize}[leftmargin=*]
    \item \textbf{身高：}突增高峰年增长10-14cm，男童$>$女童。
    \item \textbf{体重：}每年增长4-7kg，突增高峰年增长8-10kg。
\end{itemize}

\textbf{4. 生长停滞期：}青春后期开始。
\end{keybox}

\subsection{身体比例变化}

\begin{keybox}
\textbf{二、身体比例变化}

\textbf{1. 反映横向身体比例指标：}反映身体各横截面指标的相互关系。
\begin{itemize}[leftmargin=*]
    \item \imp{身高胸围指数}（胸围/身高）和\imp{肩盆宽指数}（肩宽/盆宽）：只能衡量局部（身体上部或下部），实测值重复性差，后者个体差异只在青春期显现。
    \item \imp{BMI}：既能反映身体胖瘦，又能确切反映不同群体、年代身体充实度变化。许多发达国家特别重视对BMI高百分位数变化的动态监测，将其作为制定肥胖防治策略和措施的重要依据。
\end{itemize}

\textbf{2. 反映纵向身体比例指标：}反映个体下肢长度和线性高度（身高）的比例关系。
\begin{itemize}[leftmargin=*]
    \item \imp{身高坐高指数} = 坐高/身高。
    \item \imp{下肢长指数：}下肢长指数I = (身高$-$坐高)/身高 $\times$ 100\%，下肢长指数II = (身高$-$坐高)/坐高 $\times$ 100\%。
\end{itemize}
\end{keybox}

\subsection{体型变化}

\begin{defbox}
\textbf{三、体型变化}

\textbf{1. 体型（Somatotype）：}对人体形态的总体描述和评定，是身体各部位大小比例的形态特征总和，由遗传因素决定，也受人体对环境的适应和诸多环境因素的综合影响。

\textbf{2. 分类}
\begin{itemize}[leftmargin=*]
    \item \imp{内胚层体型（Endomorphy）}：身体圆胖，消化器官发达。
    \item \imp{中胚层体型（Mesomorphy）}：身体健壮，骨骼肌肉发达。
    \item \imp{外胚层体型（Ectomorphy）}：身体瘦长，神经感觉系统发达。
\end{itemize}
\end{defbox}

% ----- 5.2 体能发育 -----
\section{儿童少年体能发育}

\subsection{概念}

\begin{keybox}
\textbf{体能（Physical Fitness）/ 体适能：}人体具备的能胜任日常工作和学习而不感到疲劳，同时有余力能充分享受休闲娱乐生活，又可应付突发紧急情况的能力。

\textbf{一、概念}

\textbf{1. 健康相关体能（Health-Related Physical Fitness）：}体能的基础部分，其目标是为身体健康和高质量，泛指人的体质，维持身体健康、提高工作、学习和生活效率所必需的基本技能。
\begin{itemize}[leftmargin=*]
    \item 指标体系：体成分、心肺耐力、柔韧性、肌力、肌耐力。
    \item 是儿童少年为维持健康、提高生活质量等需求的追求。
\end{itemize}

\textbf{2. 运动技能相关体能（Sports-Related Physical Fitness）：}建立在健康相关体能基础上，属于较高体能需求层次，目标是为比赛取胜奖项。

\textbf{运动技能相关体能六项构成要素：}
\begin{itemize}[leftmargin=*]
    \item \imp{平衡（Balance）}、\imp{协调（Coordination）}、\imp{敏捷（Agility）}
    \item \imp{速度（Speed）}、\imp{力量（Strength）}、\imp{反应时间（Reaction Time）}
\end{itemize}
\end{keybox}

\subsection{儿童少年体能发育特点}

\begin{keybox}
\textbf{二、儿童少年体能发育特点}

\textbf{1. 不均衡性：}不同体能指标在不同年龄发育速度有快有慢。

\textbf{2. 阶段性}
\begin{itemize}[leftmargin=*]
    \item 快速增长阶段：男童6-14岁，女童6-12岁。
    \item 慢速增长阶段：男童15-18岁，女童12-15岁。
    \item 恢复性生长：仅女童16-18岁。
    \item 稳定阶段：男性19-25岁，女性19-22岁。
\end{itemize}

\textbf{3. 不平衡性}

\textbf{4. 性别特征}
\end{keybox}

\subsection{儿童少年体力活动}

\begin{keybox}
\textbf{三、儿童少年体力活动}

\textbf{体力活动（Physical Activity）：}由骨骼肌收缩产生需要能量消耗的任何身体运动，包括工作、家务、旅游、游戏、娱乐等。

\textbf{体力活动与锻炼的区分：}
\begin{itemize}[leftmargin=*]
    \item 锻炼（Exercise）是体力活动的\imp{子范畴}：有计划、结构化、重复性，目的是改善或维持一种或多种体能。
    \item 体力活动涵盖范围更广（工作、家务、休闲等），锻炼是其中\imp{有目的、有组织}的那部分。
\end{itemize}

\textbf{WHO儿童青少年体力活动推荐量（5$\sim$17岁）：}
\begin{itemize}[leftmargin=*]
    \item \imp{每天累计$\geq$60分钟}中、高强度体力活动（MVPA）。
    \item 超过60分钟可提供\imp{额外的健康效益}。
    \item 每周至少\imp{3次}增强肌肉和骨骼的活动（负重/抗阻活动）。
\end{itemize}

\textbf{体力活动的五大健康作用：}
\begin{enumerate}[leftmargin=*]
    \item \imp{改善身体健康}：提高心肺功能和肌肉骨骼健康。
    \item \imp{改善心血管代谢健康}：改善血压、血脂异常、血糖控制和胰岛素抵抗（以有氧运动为主）。
    \item \imp{提高骨矿物质含量和骨密度}：预防成年后骨质疏松症和相关骨折。
    \item \imp{改善认知功能}：对学业成绩、记忆力和执行功能有积极影响。
    \item \imp{提高心理健康水平}：降低抑郁症状和抑郁发生风险。
\end{enumerate}
\end{keybox}

% ----- 5.3 体成分发育 -----
\section{儿童少年体成分发育}

\subsection{体成分模型}

\begin{defbox}
\textbf{身体成分（Body Composition）/ 体成分：}人体总重量中不同身体成分的构成比例，属化学生长（Chemical Growth）范畴。

\textbf{一、体成分模型}

\textbf{1. 二成分模型：}体脂重（FM）+ 去脂体重（FFM），主要使用。

\textbf{2. 多成分模型：}二成分模型的细化。
\end{defbox}

\subsection{体成分发育的性别特征}

\begin{defbox}
\textbf{二、体成分发育的性别特征}

\textbf{1. 体成分发育的性别差异}
\begin{itemize}[leftmargin=*]
    \item 生命早期，男女体脂量接近。
    \item 儿童期，男、女体脂均下降。
    \item 青春期，女童变化以体脂增加为主，男童体脂有下降。
    \item 体脂分布的性别差异在青春期前即有表现；男童上半身的皮下和内脏脂肪蓄积，女童皮下脂肪蓄积在臀部。
\end{itemize}

\textbf{2. 瘦体重发育的性别特征}
\begin{itemize}[leftmargin=*]
    \item 生命早期开始，男孩就显出更高的瘦体重，故其体重$>$同年龄女孩。
    \item 儿童期，男孩仍较女孩有较多瘦体重增长。
    \item 进入青春期后，男孩主要表现为瘦体重增加，女孩的瘦体重增量仅为男孩的一半。
    \item 随着近年来性早熟（女性更多）现象增加，瘦体重的性别差异出现时间明显提前。
\end{itemize}
\end{defbox}

\subsection{体成分发育的年龄特征}

\begin{keybox}
\textbf{二、体成分发育的年龄特征}

\textbf{1. 体脂发育的年龄变化：}
\begin{itemize}[leftmargin=*]
    \item 男、女童在青春期突增前体脂率（\%BF）均上升。
    \item 男孩：增长高峰在11岁，其后下降，青春后期再度持续上升，18岁平均值约17.0\%。
    \item 女孩：整个青春期\%BF都上升，突增高峰也不下降，18岁时平均为27.8\%。
\end{itemize}

\textbf{2. 瘦体重随年龄的变化：}
\begin{itemize}[leftmargin=*]
    \item 青春期前，瘦体重随年龄增长而增加，女孩的瘦体重增长不如男孩显著，性别差异随年龄增长而扩大。
    \item 青春期，瘦体重的变化主要是含水量逐年减少，女孩总体水量仍高于男孩。
\end{itemize}

\textbf{3. 其他体成分的年龄特征：}
\begin{itemize}[leftmargin=*]
    \item 细胞外水随年龄增长而逐渐浓缩，儿童肌肉中水分比成人多，伴随年龄增长而相对减少。10--18岁时男、女孩总体水出现差异。
    \item 男、女童体内钙含量的改变与生长突增关系密切，男童总体钙的均值在13--14岁增加达35\%；女童11--12岁增加最多（约40\%）。
    \item 女童骨矿盐的沉积量大约1/3发生在青春期开始后的3--4年，而男童直到15--18岁仍有骨矿盐的沉积。
\end{itemize}
\end{keybox}

\subsection{体成分发育的种族差异}

\begin{defbox}
\textbf{三、体成分发育的种族差异}

\textbf{1. 体脂的种族差异：}
\begin{itemize}[leftmargin=*]
    \item 体型和青春期发育状态都会受到种族影响，因此导致体脂的种族差异。
    \item 移民流行病学研究显示：西班牙裔白人体脂率最高；非裔黑人的内脏脂肪含量比欧裔白人低。
\end{itemize}

\textbf{2. 瘦体重的种族差异：}
不同种族的瘦体重差异较大。研究表明，无论男女，非洲裔儿童青少年的去脂体质量指数（FFMI）趋势均高于白人和拉丁裔人。
\end{defbox}

\subsection{儿童期体成分发育的健康效应}

\begin{keybox}
\textbf{四、儿童期体成分发育的健康效应}

\textbf{（一）体成分对身体素质与机能的影响}

\textbf{1. 对身体素质的影响：}
\begin{itemize}[leftmargin=*]
    \item 身体素质是个体体质强弱的外在表现，无论是速度、爆发力，还是有氧耐力都需要一定的体脂肪和去脂体重基础。
    \item 研究表明，去脂体质量指数（FFMI）较高、体脂率（FM\%）和体脂指数（FMI）较低的学生在肌肉爆发力、肌肉耐力、柔韧性、有氧能力、无氧能力方面有更好表现，且这种趋势在男生中更为明显。
\end{itemize}

\textbf{2. 对于身体机能的影响：}
\begin{itemize}[leftmargin=*]
    \item 身体机能涉及人体各组织、器官和系统的生理活动和整体表现，常用指标包括心率、血压和肺活量等，体脂肪率对这些机能指标有显著影响。
    \item 研究表明，7岁儿童的FM\%与血压呈正相关，但与肺活量呈显著负相关；这一趋势在7--12岁的小学生中也存在，且显示出一定的剂量反应关系。
    \item 青少年的血压偏高检出率随着体脂肪率的上升而增加，尤其在男生中这种关联更为明显。
\end{itemize}

\textbf{（二）体成分对儿童发育的影响}

\textbf{1. 对青春期发育的影响：}
\begin{itemize}[leftmargin=*]
    \item 儿童的体脂量和体脂率对其青春期的启动具有显著影响。
    \item 研究表明，女童青春期提前发动组的体脂率和瘦体重均高于正常发动组，后者又高于发动延迟组，体脂率过高与女性青春期的提前发动密切相关，且存在显著的剂量反应关系。
    \item 男性瘦体重的增加同样与青春期提前发动有关，表现出提前组、正常组、延迟组的递减趋势。
    \item 无论男生还是女生，高体脂率均可能引起青春期的提前启动。
\end{itemize}

\textbf{2. 对骨骼发育的影响：}
\begin{itemize}[leftmargin=*]
    \item 研究发现，在8--9岁的小学生中，骨骼肌含量与骨龄发育呈负相关，而体脂肪含量和身体水分含量则与骨龄发育正相关。
    \item 超重和肥胖的男女生骨龄发育提前的概率分别是正常体重儿童的2.148倍和2.932倍。
    \item 肌肉量的增加可能促进骨骼的健康发育，但过高的体脂肪可能对骨发育产生不利影响。
\end{itemize}

\textbf{（三）身体成分对慢性非传染性疾病发病的影响}

\textbf{1. 与心血管健康：}
\begin{itemize}[leftmargin=*]
    \item 儿童期体成分异常，尤其是体脂率偏高，与成年后患心血管疾病的风险增加有关。
    \item 青春期学生体脂率与血压偏高正相关，尤其在男生中更为明显。
    \item 儿童期身体成分分析能更客观反映脂肪含量和分布，对预防心血管风险具有指导意义。
\end{itemize}

\textbf{2. 与内分泌健康：}
\begin{itemize}[leftmargin=*]
    \item 儿童躯干部位的体脂百分比是胰岛素抵抗的独立危险因素。
    \item 内脏脂肪（VAT）与肥胖青少年胰岛素抵抗的早期发展密切相关，随着VAT的增加，腹部皮下脂肪减少，与肝脂肪变性、胰岛素抵抗和代谢综合征风险增加有关。
    \item 肥胖和超重是儿童罹患代谢综合征的重要风险因素，体脂分布与儿童代谢综合征的关联紧密，尤其是腹部肥胖，明显提升发病风险。
\end{itemize}
\end{keybox}

\newpage
